\documentclass[
	12pt, %Schriftgröße
	a4paper, %Papierformat
	oneside, %einseitiger Druck, für zweiseitigen Druck weglassen
%	openany, %neue Kapitel dürfen auch auf linker Seite beginnen (beidseitiger Druck), sonst openright (default)
	chapterprefix, %Kapitelüberschriften zu Beginn eines Kapitels anzeigen
	listof=totoc, %Verzeichnisse (Abbildungen/Tabellen) im Inhaltsverzeichnis aufführen
%	draft, %Entwurfsmodus: lässt Links und Abbildungen weg, zeigt übervolle Zeilen an, wird schneller gesetzt
]{scrbook}

\newcommand*{\universitaet}[1]{\def\universitaet{#1}}
\newcommand*{\fach}[1]{\def\fach{#1}}
\newcommand*{\arbeitsart}[1]{\def\arbeitsart{#1}}
\newcommand*{\titel}[1]{\def\titel{#1}}
\newcommand*{\autor}[1]{\def\autor{#1}}
\newcommand*{\gutachter}[1]{\def\gutachter{#1}}
\newcommand*{\zweitgutachter}[1]{\def\zweitgutachter{#1}}
\newcommand*{\abteilung}[1]{\def\abteilung{#1}}
\newcommand*{\department}[1]{\def\department{#1}}
\newcommand*{\ort}[1]{\def\ort{#1}}
\newcommand*{\datum}[1]{\def\datum{#1}}

%---------------------------------
%---Sprache, Schrift & Encoding---
%---------------------------------

\usepackage[ngerman]{babel} %neue deutsche Rechtschreibung
\usepackage[utf8]{inputenc} %alternativ: ansinew (windows), latin1, latin9 (Linux/UNIX/Mac)

%\usepackage[T1]{fontenc} %einkommentieren, wenn Buchstaben/Schriftbild schief aussehen!
%\usepackage{lmodern} %andere Schrift, sieht evtl. besser aus

%------------
%---Layout---
%------------

\usepackage[
	left=3cm, %linker Seitenrand
	right=3cm, %rechter Seitenrand
	top=2.5cm, %oberer Seitenrand
	bottom=2.5cm, %unterer Seitenrand
]{geometry}
\usepackage[onehalfspacing]{setspace} %1,5-facher Zeilenabstand

%\usepackage{indentfirst} %rückt auch den ersten Absatz eines (Unter)abschnitts ein, so wie alle anderen Absätze auch

\usepackage[
	headsepline, %Kopfzeile durch Querlinie abheben
%	footsepline, %Fußzeile durch Querlinie abheben
	draft=false, %Kopf-/Fußzeile auch im draft-Modus (siehe documentclass) anzeigen
]{scrlayer-scrpage}
\clearpairofpagestyles
\ihead[]{\headmark} %Kapitel-/Abschnittsname innen in der Kopfzeile, außer zu Beginn neuer Kapitel
%\chead[]{}
\ohead*[]{\pagemark} %Seitenzahl außen in der Kopfzeile, außer zu Beginn neuer Kapitel
%\ifoot[]{}
%\cfoot[]{}
\ofoot*[\pagemark]{} %Seitenzahl außen in die Fußzeile, aber nur zu Beginn neuer Kapitel

\renewcommand*\chapterformat{\chapapp~\thechapter} %Kapitelüberschriften in der Form Kapitel~<Nummer>, siehe dazu chapterprefix in der documentclass

%-----------
%---Mathe---
%-----------

\usepackage{amsmath}
\usepackage{amsfonts}
\usepackage{amssymb}

%---------------
%---Sonstiges---
%---------------

\usepackage[sort&compress]{natbib} %ermöglicht die Verwendung von \citet{} und \citep{}
\bibliographystyle{dinat} %Zitierstil: andere auch möglich, z.B. alphadin, plain, ...

\usepackage{caption}
\usepackage{subcaption}
\usepackage{graphicx}
\usepackage{color}
\usepackage{transparent}
\usepackage{enumerate}
\usepackage{longtable}
\usepackage{float}
\restylefloat{figure}
\graphicspath{{Abbildungen/}}

\usepackage[
	textsize=tiny, %kleinere Schrift als die default-Größe in den Notizen, damit mehr reinpasst
]{todonotes} %Notizen im Text mit

%------------------
%---Links & URLs---
%------------------

%hyperref möglichst als letztes package laden (es gibt aber Ausnahmen)!

\usepackage[
	urlcolor=black, %Textfarbe von Links schwarz statt blau (default)
]{hyperref} %für Referenzen, Verweise, Links, ...

%-----------------------------
%---Variablen zum ausfüllen---
%-----------------------------

\universitaet{Carl von Ossietzky Universität Oldenburg}
\fach{Wirtschaftsinformatik} %Informatik, ...
\arbeitsart{Bachelorarbeit} %Masterarbeit, Diplomarbeit, ...
\titel{Vergleich von verschiedenen Deep Reinforcement Learning Agenten am Beispiel des Videospiels Snake}
\autor{Lorenz Mumm}
\gutachter{Apl. Prof. Dr. Jürgen Sauer}
\zweitgutachter{M. Sc. Julius Möller}
\abteilung{Abteilung Systemanalyse und -optimierung}
\department{Department für Informatik}
\ort{Oldenburg}
\datum{\today} %Tag der Abgabe BEIM PRÜFUNGSAMT

%----------------
%---Titelseite---
%----------------

\begin{document}

\frontmatter %römische Seitennummerierung für den Anfang



\makeatletter
\begin{titlepage}
\centering
\begin{figure}[htb]
  \centering
  \begin{subfigure}{0.5\textwidth}
    \centering
    \includegraphics[scale=1.8]{Abbildungen/uni_ol_logo.png}
    \end{subfigure}%
    \begin{subfigure}{0.5\textwidth}
      \centering
      \includegraphics[scale=0.35]{Abbildungen/abteilungslogo.jpg}
    \end{subfigure}
\end{figure}

{\scshape\Large\universitaet\par}
\vspace{1cm}
{\scshape\normalsize\fach\par}
{\scshape\normalsize\arbeitsart\par}

\vfill

\hrule
\vspace{1cm}
{\LARGE\bfseries\titel\par}
\vspace{1cm}
\hrule

\vfill

\begin{minipage}[t]{0.4\textwidth}
  \begin{flushleft}
    {\normalsize\itshape Autor:\par}
    {\normalsize\autor}
  \end{flushleft}
\end{minipage}%
\begin{minipage}[t]{0.4\textwidth}
  \begin{flushright}
    {\normalsize\itshape Erstgutachter:\par}
    {\normalsize\gutachter}
  \end{flushright}
  \begin{flushright}
    {\normalsize\itshape Zweitgutachter:\par}
    {\normalsize\zweitgutachter}
  \end{flushright} 
\end{minipage}

\vfill

{\normalsize\abteilung\par}
{\normalsize\department\par}
\vspace{1cm}
{\normalsize\ort{}, \datum\par}
\end{titlepage}
\makeatother

%---------------
%---Abstracts---
%---------------

%einkommentieren, falls benötigt
%\chapter*{Zusammenfassung}
Auf Deutsch.
%\chapter*{Abstract}
In English.

%-------------------
%---Verzeichnisse---
%-------------------

\tableofcontents %Inhaltsverzeichnis erzeugen
\listoffigures %Abbildungsverzeichnis erzeugen
\listoftables %Tabellenverzeichnis erzeugen

\chapter{Abkürzungsverzeichnis}
\begin{center}
\begin{tabular}{rl}
\textbf{KI} & \textbf{K}ünstliche \textbf{I}entelligenz\\
\textbf{RL} & \textbf{R}einforcement \textbf{L}earning \\
\textbf{DQN} & \textbf{D}eep \textbf{Q}-\textbf{N}etwork\\
\textbf{DQL} & \textbf{D}eep \textbf{Q}-\textbf{L}earning\\
\textbf{DDQN} & \textbf{D}ouble \textbf{D}eep \textbf{Q}-\textbf{N}etwork\\
\textbf{PPO} & \textbf{P}roximal \textbf{P}olicy \textbf{O}ptimization\\
\textbf{SAC} & \textbf{S}oft \textbf{A}ctor \textbf{C}ritic\\
\textbf{A2C} & Advantage \textbf{A}ctor \textbf{C}ritic\\
\textbf{Env} & \textbf{Env}ironment \\
\textbf{Obs} & \textbf{Obs}ervation \\
\end{tabular}
\end{center}
  %auskommentieren, falls nicht benötigt

%---------------
%---Hauptteil---
%---------------

\mainmatter %arabische Seitennummerierung und nummerierte Kapitel ab dem Hauptteil

%die Kapitel der Arbeit in einzelnen Dateien anlegen und mit \input{} importieren
\chapter{Einleitung}\label{sec:beispiel}
Das Maschine Learning ist weltweit auf dem Vormarsch und große Unternehmen investieren riesige Beträge, um mit KI basierten Lösungen größere Effizienz zu erreichen. Auch der Bereich des Reinforcement Learning gerät dabei immer mehr in das Blickfeld von der Weltöffentlichkeit. Besonders im Gaming Bereich hat das Reinforcement Learning schon beeindruckende Resultate erbringen können, wie z.B. die KI AlphaGO, welche den amtierenden Weltmeister Lee Sedol im Spiel GO besiegt hat \cite{UAV}. In Anlehnung an die vielen neuen Reinforcement Learning Möglichkeiten, die in der letzten Zeit entwickelt wurden und vor dem Hintergrund der immer größer werdenden Beliebtheit von KI basierten Lösungsstrategien, soll es in dieser Bachelorarbeit darum gehen, einzelnen Reinforcement Learning Agenten, mittels statistischer Methoden, genauer zu untersuchen und den optimalen Agenten für ein entsprechendes Problem zu bestimmen.

\section{Motivation} \label{sec:Motivation}
In den letzten Jahren erregte das Reinforcement Learning eine immer größere Aufmerksamkeit. Siege über die amtierenden Weltmeister in den Spielen GO oder Schach führten zu einer zumehmenden Beliebtheit des RL. Neue Verfahren, wie z.B. der Deep Q-Network Algorithmus auf dem Jahr 2015 \cite{DQN}, der Proximal Policy Optimization (PPO) aus dem Jahr 2017 \cite{PPO} oder der Soft Actor Critic (SAC) aus dem Jahr 2018 \cite{SAC}, haben ihr Übriges getan, um das RL auch in anderen Bereichen weiter anzusiedeln, wie z.B. der Finanzwelt, im autonomen Fahren, in der Steuerung von Robotern, in der Navigation im Web oder als Chatbot \cite{DRL_Lapan}.\\
Durch die jedoch große Menge an RL-Verfahren gerät man zunehmend in die problematische Situaion, sich für einen diskreten RL Ansatz zu entscheiden. Weiter erschwert wird dieser Auswahlprozess noch durch die Tatsache, dass die einzelnen Agenten jeweils untereinander große Unterschiede aufweisen. Auch existieren häufig mehrere Ausprägungen eines RL-Verfahrens, wie z.B. der Deep Q-Network Algorithmus (DQN) und der Double Deep Q-Network Algorithmus (DDQN). Die Wahl des passenden Agenten kann großen Einfluss auf die Performance und andere Bewertungskriterien haben, (\cite{Exploration_of_Reinforcement_Learning_to_SNAKE}), deshalb soll in dieser Ausarbeitung eine Vorgehen, welches auf einem Vergleich beruht, entwickelt werden, dass den optimalen Agenten für eine entsprechendes Problem bestimmt.\\
\\Eine potenzielle Umgebung, in welcher Agenten getestet und verglichen werden können, ist das Spiel Snake.
Mit der Wahl dieses Spieles ist zusätzlich zu dem oben erwähnten Mehrwert noch ein weiterer in Erscheinung getreten.\\ 
So interpretieren neue Forschungsansätze das Spiel Snake als ein dynamisches Pathfinding Problem. Mit dieser Art von Problemen sind auch unbemannte Drohne (UAV Drohnen) konfrontiert, welche beispielsweise Menschen in komplexen Katastrophensituationen, wie z.B. auf explodierten Rohölbohrplattformen oder Raffinerien, finden und retten sollen. Auch kann das Liefen von wichtigen Gütern, beispielsweise medizinischer Natur, in solche Gebiete kann durch die Forschung am Spiel Snake möglich gemacht werden.(\cite{UAV})\\
Somit kann der RL-Vergleich am Spiel Snake möglicherweise helfen, Menschen zu retten.

\section{Zielsetzung}
Ziel der Arbeit ist es für die Umgebung Sanke einen optimalen Agenten zu bestimmen, welcher spezifische Anforderungen erfüllen soll.
Damit ergib sich die folgende Forschungsfrage:\\
Wie kann an einem Beispiel des Spiels Snake für eine nicht-triviale gering-dimensionale Umgebung ein möglichst optimaler RL Agent ermittelt werden?\\
\\Basierend auf der Forschungsfrage ergibt sich ein Mehrwert für die Wissenschaft. Durch die Abnahme des Entscheidungsfindungsprozesses müssen Forscherinnen und Forscher wie auch Anwenderinnen und Anwender von RL-Verfahren nicht mehr unreflektiert irgendeinen RL Agenten auswählen, sondern können auf Grundlage der Methodik und der daraus hervorgehenden Daten den passenden Agenten bestimmen.



\chapter{Grundlagen}\label{sec:Grundlagen}

\section{Game of Snake}
Snake (englisch für Schlange) zählt zu den meist bekannten Computerspielen unserer Zeit. Es zeichnet sich durch sein simples und einfach zuverstehendes Spielprinzip aus.
In seiner ursprünglichen Form ist Snake als ein zweidimensionales rechteckiges Feld. Dieses Beschreibe das komplette Spielfeld, in welchem man sich als Snake bewegt. Häufig wird diese als einfacher grüner Punkt (Quadrat) dargestellt. Dieser stellt den Kopf der Snake dar. Neben dem Kopf der Snake befindet sich auf dem Spielfeld auch noch der sogenannte Apfel. Dieser wird häufig als roter Punkt (Quadrat) dargestellt.
\begin{figure}[H]
	\centering
	\def\svgscale{0.90}
	\input{Abbildungen/Game_of_Snake.pdf_tex}
	\caption[Game of Snake]{Game of Snake - Abbildung eines Snake Spiels in welchem der Apfel druch das rote und der Snake Kopf durch das dunkelgrüne Quadrat dargestellt wird. Die hellgrünen Quadrate stellen den Schwanz der Snake dar.}
	\label{fig:Game_of_Snake}
\end{figure}
Ziel der Snake ist es nun Äpfel zu fressen. Dies geschied, wenn der Kopf der Snake auf das Feld des Apfels läuft. Danach verschwindet der Apfel und ein neuer erscheint in einem noch freies Feld). Außerdem wächst, durch das Essen des Apfels, die Snake um ein Schwanzgglied. Diese Glieder folgen dabei in ihren Bewegungen den vorrangegangenen Schwanzgleid bishin zum Kopf. Dem Spieler ist es nur möglich den Kop der Snake zu steuern.\\
Bisher so weit so gut, doch es existiert noch einen Hacken an der Geschichte. Der Snake ist es nicht erlaubt die Wände oder sich selbst zu berühren, geschied dies, im Laufe des Spiels, trotzdem endet dieses sofort. Diese Einschränkung führt zu einem Ansteigen der Komplexität gegen Ende des Spiels. Ein Spiel gilt als gewonnen, wenn es der Snake gelungen ist, das komplette Spielfeld auszufüllen.

\section{Reinforcement Learning}
Das Reinforcement Learning (Bestärkendes Lernen) ist einer der drei großen Teilbereiche, die das Machine Learning zu bieten hat. Beben dem Reinforcement Learning zählen das supervised Learning (Überwachtes Lernen) und das unsupervised Learning (unüberwachtes Lernen) ebenfalls noch zum Machine Learning.\\
Einordnen lässt sich das Reinforcement Learning (RL) als Supervised Learning mit fehlenden Labels. Viele Aspekte des (SL), wie z.B. neuronale Netze zum Approximation einer Lösungsfunktion, Gradientenabstiegsverfahren und Backpropergation zur Erlernung von Datenrepräsentationen, werden auch im RL verwendet.\\
Auf Menschen wirkt das RL, im Vergleich zu den anderen Disziplinen des Machine Learnings, am nachvollziehbarsten. Dies liegt an der Lernstrategie die sich dieses Verfahren zu Nutze macht. Beim RL wird anhand eines 
\grqq trial-and-error\glqq{} Verfarens gelernt. Ein gutes Beispiel für eine solche Art des Lernens ist die Erziehung eines Kindes. Wenn eben dieses Kind etwas gutes tut, dann wird es belohnt. Angetrieben von der Belohnung, versucht das Kind dieses Verhalten fortzusetzen. Entsprechend wird das Kind bestraft, wenn es etwas schlechtes tut. Schlechtes Verhalten kommt weniger häufig zum Vorschein, um Bestrafungen zu vermeiden.\\
Beim RL funktioniert es genau so. Das ist auch der Grund dafür, dass viele der Aufgaben des RL dem menschlichen Arbeitsspektrum sehr nahe sind. So wird das RL beispielsweise im Finanzhandel eingesetzt. Was früher Börsianer getan haben, erledigt heute der RL Agent. Auch im Bereich der Robotik ist das RL auf dem vormarsch. Wo früher noch komplexe Bewegungsabfolgen eines Roboterarms mühevoll programmiert werden mussten, da können wir heute bereits Roboter mit RL Agenten ausstatten, welche selbstständig die Bewegungsabfolgen meisten \cite{DRL}.

\subsection{Klassifikation von RL-Verfahren}
Bevor ein tiefer technischer Einblick in das RL vorgenommen werden sollte, ist es ratsam sich erst einmal mit den einzelnen Klassifikationen des RL zu beschäftigen, um ein tiefergehendes Verständnis für die einzelnen RL Verfahren zu entwickeln. Alle RL Verfahren lassen sich, unter gewissen Gesichtspunkten, in Klassen einordnen, welche Aufschluss über die Implementierung, den Entscheidungsfindungsprozess und die Dateneffizens geben. 
Natürlich existieren viele Möglichkeiten RL-Verfahren zu klassifizieren aber vorerst beschränken wir uns auf die folgenen drei.

\subsubsection{Model-free und Model-based}
Die Unterscheidung in model-free (modelfrei) und in model-based (modelbasiert) gibt Aufschluss darüber, ob der Agent fähig ist, ein Modell des Zustandekommens der Belohnungen (Reward) zu entwickeln.\\
Model-free RL Verfahren sind nicht in der Lage das Zustandekommen der Belohnung vorherzusagen, vielmehr ordnen sie einfach die Beobachtung einer Aktion zu. Es werden keine zukünftigen Beobachtungen und/oder Belohnungen extrapoliert.\\
Ganz anderes sieht es da bei den model-based RL Verfahren aus. Diese versuchen eine oder mehrere zukünftige Beobachtungen und/oder Belohnungen zu ermitteln, um die beste Aktionsabfolge zu bestimmen. Dem Model-based RL Verfahren liegt also ein planungsprozess der nächsten Züge zugrunde \cite{Sutton1998}.\\
\\Beide Verfahrensklassen haben Vor- und Nachteile, so sind model-based Verfahren häufig in deterministischen Environments mit simplen Aufbau und strengen Regeln anfindbar. Die Bestimmung von Observations und/oder Rewards bei größerern Environments. wäre viel zu komplex und zu resourcenbindend.
\\Model-Free Algorithmen haben dagegen den Vorteil, dass das Trainieren leichter ist, aufgrund des wegfallenden Aufwandes, welcher mit der Bestimmung zukünftiger Observations und/oder Rewards einhergeht. Sie preformen zudem in großen Environments besser als model-based RL Verfahren. Des Weiteren sind model-free RL Verfahren universiell einsetzbar, im Gegensatz zu model-based Verfahren, welche ein Modell des Environment für das Planen benötigen  \cite{DRL, Sutton1998}.

\subsubsection{Policy-based und value-based Verfahren}
\label{RL_policy_value}
Die Einordnung in Policy- und wertebasierte Verfahren gibt Aufschluss über den Entscheidungsfindungsprozess des Verfahrens.
Agenten, welche policy-based arbeiten, versuchen unmittelbar die berechnete Policy umzusetzen und sie zu optimieren. Policy-based RL Verfahren besitzen dafür meist ein eigenes Network (Policy-Network), welches die Policy bestimmt.
Gewöhnlich wird diese als eine Wahrscheinlichkeitsverteilung über alle Actions repräsentiert. Jede Action erhält damit einen Wert zwischen Null und Eins, welcher Aufschluss über die Qualität der Action im momentanen Zustand des Env. liefert.\\
Anders als bei den policy-based wird bei den value-based Verfahren nicht mit Wahrscheinlichkeiten gearbeitet. Die Policy und damit die Entscheidungsfindung, wird indirekt über das Bestimmen aller Values ermittelt. Es wird daher immer die Action gewählt, welche zum dem State führt, der über den größten Value verfügt \cite{DRL}.

\subsubsection{On-Policy und Off-Policy Verfahren}
Eine Klassifikation in On-Policy und Off-Policy Verfahren hingegen, gibt Aufschluss über die Dateneffizenz. Einfach formuliert sind RL Verfahren, welche Off-Policy sind, in der Lage von älteren Daten zu lernen. Diese können vom Menschen oder von Agenten generiert worden sein. Es spielt keine Rolle, mit welcher Entscheidungsfindungsqualität die Daten erhoben worden sind. Sie können zu Begin, in der Mitte oder zum Ende des Lernprozesses ermittelt worden sein. Die Aktualität der Daten spielt daher keine Rolle für Off-Policy-Verfahren.\\
On-Policy-Verfahren sind dagegen sehr wohl abhängig von aktuellen Daten, da sie versuchen die Policy indirekt oder direkt zu optimieren.\\
\\Auch hier besitzen beide Klassen ihre Vor- und Nachteile. So können beispielsweise Off-Policy Verfahren mit älteren Daten immer noch trainiert werden. Dies macht Off-Policy RL Verfahren dateneffizienter als On-Policy Verfahren. Meist ist es jedoch so, dass diese Art von Verfahren langsamer konvergieren.\\
On-Policy Verfahren konvergieren dagegen meist schneller. Sie benötigen aber dafür auch mehr Daten aus dem Environment, dessen Beschaffung aufwendig und teuer sein könnte. Die dateneffizienz nimmt ab \cite{DRL}.

\subsection{Vokabular}
Um jedoch das RL tiefergehend zu verstehen, ist es erforderlich die gängigen Begrifflichkeiten zu erlernen und deren Bedeutung zu verstehen.

\subsubsection{Agent}\label{sec:Agent}
Im Zusammenhang mit dem RL hört man gewöhnlich von Agenten. Sie sind die zentralen Kapselungsobjekte, welche die eigentlichen Algorithmen, wie z.B. den Algorithmus des Q-Learning oder eines Proximal Policy Optimization, in eine feste Schnittstelle einbinden. Häufig sind auch Methoden in der Schnittstelle vorgegeben, wie z.B. die choose\_action oder die learn Methode. Je nach Implementierung können die Methoden gelegentlich auch andere aber ersichtliche Namen tragen, wie z.B. die choose\_action Methode, welche auch häufig unter dem Namen \glqq act\grqq{} anzutreffen ist.\\
Bei dem Agenten handelt es sich gewöhnlich um die einzige Instanz, welche mit dem Environment (der Umgebung) interagiert. Zu dieser Interaktionen zählen das Entgegennehmen von Observations und Rewards, wie auch das Tätigen von Actions \cite{DRL}.

\subsubsection{Environment}\label{sec:Environment}
Das Environment (Env.) bzw. die  Umgebung ist vollständig außerhalb des Agenten angesiedelt. Es spannt das zu manipulierende Umfeld auf, inwelchem der Agent Interaktionen tätigen kann. An ein Environment werden unter anderem verschiedene Ansprüche gestellt, damit ein RL Agent mit ihm in Interaktion treten kann. Zu diesen Ansprüchen gehört unter anderem die Fähigkeit Observations und Rewards zu liefen, Actions zu verarbeiten \cite{DRL}.
Actions, welche im momentanen State des Env. nicht gestattet sind, müssen entsprechend behandelt werden. 
Dies wirft die Frage auf, ob es in dem Env. einen terminalen Zustand (häufig done oder terminal genannt) geben soll. Existiert ein solcher terminaler Zustand, so muss eine Routine für den Reset (Neustart) des Env. implementiert sein.

\subsubsection{Action}\label{sec:Action}
Die Actions bzw. die Aktionen sind einer der drei Datenübermittlungswege. Bei ihnen handelt es sich um Handlungen welche im Env. ausgeführt werden. Actions können z.B. erlaubte Züge in Spielen, das Abbiegen im Autonomenfahren oder das Ausfüllen eines Antragen sein. Es wird ersichtlich, dass die Actions, welche ein RL Agent ausführen kann, prinzipiell nicht in der Komplexität beschrängt sind. Dennoch ist es hier gängige Praxis geworden, dass arbeitsteilig vorgegangen werden soll, da der Agent ansonsten zu viele Ressourcen in Anspruch nehmen würde.\\
Im RL wird zwischen diskreten und stetigen Aktionsraum unterschieden. Der diskrete Aktionsraum umfassen eine endliche Menge an sich gegenseitig ausschließenden Actions. Beispielhaft dafür wäre das Gehen an einer T-Kreuzung. Der Agent kann entweder Links, Rechts oder zurück gehen. Es ist ihm aber nicht möglich ein bisschen Rechts und viel Links zu gehen.
Anders sicht das beim stetigen Aktionsraum aus. Dieser zeichnet sich durch stetige Werte aus. Hier ist das Steuern eines Autos beispielhaft. Um wie viel Grad muss der Agent das Steuer drehen, damit das Fahrzeug auf der Spur bleibt \cite{DRL}?

\subsubsection{Observation}\label{sec:Observation}
Die Oberservation (Obs.) bzw. die Beobachtung ist ein weiterer der drei Daten-übermittlungswege, welche den Agenten mit dem Environment verbindet. Mathematisch, handelt es sich bei der Obs. um einen oder mehere Vektoren bzw. Matrizen. Die Obs. hat einen immensen Einfluss auf den Erfolg des Agenten und sollte daher klug gewählt werden.
Je nach Anwendungsbereich fällt die Obs. sehr unterschiedlich aus. In der Finanzwelt könnte diese z.B. die neusten Börsenkurse einer oder mehrerer Aktien beinhalten oder in der Welt der Spiele könnten diese die aktuelle erreichte Punktezahl wiedergeben \cite{DRL}. 
Es hat sch als Faustregel herausgestellt, dass man sich bei dem Designing der Obs. auf das wesentliche konzentieren sollte. Unnötige Informationen können den die Effizenz des Lernen mindern und den Resourcenverbrauch zudem steigen lassen.

\subsubsection{Reward}\label{sec:Reward}
Der Reward bzw. die Belohnung ist der letzte Datenübertragungsweg. Er ist neben der Action und der Obs. eines der wichtigsten Elemente des RL und von besonderer Bedeutung für den Lernerfolg.
Bei dem Reward handelt es sich um eine einfache skalare Zahl, welche vom Env. übermittelt wird. Sie gibt an, wie gut oder schelcht eine ausgeführte Action im Env. war \cite{DRL}. Um eine solche Einschätzung zu tätigen, ist es nötig eine Bewertungsfunktion zu Implementieren, welche den Reward bestimmt. 
Bei der Modellierung des Rewards kommt es voralledem darauf an, in welchen Zeitabständen dieser an den Agenten gesendet wird (sekündlich, minütlich, nur ein Mal). Aus Bequemlichkeitsgründen ist es jedoch gängige Praxis geworden, dass der Reward in fest definierten Zeitabständen erhoben und üermittelt wird \cite{DRL}.
Je nach Implementierung hat dies große Auswirkungen auf das zu erwartende Lernverhalten.

\subsubsection{State}\label{sec:State}
Der State bzw. der Zustand ist eine Widerspieglung der zum Zeitpunkt $t$ vorherschenden Situation im Environments. 
Der State wird von der Obs. repräsentiert. Häufig findet der Begriff des States in diversen Implementierungen, wie auch in vielen Ausarbeitung zum Themengebiet des RL Anwendung \cite{Sutton1998}.

\subsection{Funktionsweise}
\begin{figure}[H]
	\centering
	\def\svgscale{1.0}
	\input{Abbildungen/Reinforcement_Learning.pdf_tex}
	\caption[Reinforcement Learning]{Reinforcement Learning schematische Funktionsweise - Der Agent erhält einen State $S_{t}$ und falls $t \mathrel{{!}{=}} 0$ einen Reward $R_{t}$. Daraufhin wird vom Agenten eine Action $A_{t}$ ermittelt, welche im Environment ausgeführt wird. Das Environment übermittelt den neu enstandenen State $S_{t+1}$ und Reward $R_{t+1}$ an den Agenten. Diese Prozedue wird wiederholt \cite{Sutton1998}.}
	\label{fig:Reinforcement Learning}
\end{figure}
Zu Beginn wird dem Agenten vom Environment der initialer State übermittelt. Auf Grundlage dieses Stat $S_{t}$ wobei $t=0$ ist, welcher inhaltlich aus der zuvor besprochenden Obs. \ref{sec:Observation} besteht, wird im Agenten ein Entscheidungsfindungsprozess angestoßen. Es wird eine Action $A_{t}$ ermittelt, welcher der Agent an das Environment weiterleitet. Die vom Agenten ausgewählte Action $A_{t}$ wird num im Env. ausgeführt. Dabei kann der Agent selbstständig das Env. manipuliert oder er kann die Action an das Env. weiterleitet. Es kümmert sich dann  selbst um die Durchführung der Action.\\
Das manipulierte Environment befindet sich nun im neuen State $S_{t+1}$, welcher an den Agenten weitergeleitet wird. Des Weiteren wird noch einen Reward $R_{t+1}$, welcher vom Env. bestimmt wurde, an den Agenten übermittelt.\\ 
Mit dem neuen State $S_{t+1}$, kann der Agent wieder eine Action $A_{t+1}$ bestimmen, die ausgeführt wird. Daraufhin werden wieder der neue State $S_{t+2}$ und Reward $R_{t+2}$ ermittelt und übertragen usw. Der Zyklus beginnt von neuem \cite{Sutton1998}.

\section{Proximal Policy Optimization}
Der Proximal Policy Optimization Algorithmus oder auch PPO abgekürzt wurde von den Open-AI-Team entwickelt. Im Jahr 2017 erschien das gleichnamige Paper, welches von John Schulman et al. veröffentlicht wurde. In diesem werden die besonderen Funktionsweise genauer erläutert wird \cite{PPO}.
\subsection{Actor-Critic Modell}
Der PPO Algorithmus ist ein policy-based RL-Verfahren, welches, im Vergleich mit anderen gleichen Verfahren, einige Verbesserungen aufweist. Er ist eine Weiterentwicklung des Actor-Critic-Verfahrens und basiert intern auf zwei NN \cite{DRL}, dem sogenanntes Actor-Network bzw. Policy-Network und das Critic-Network bzw. Value-Network. Beide NN können aus mehreren Schichten bestehen, jedoch sind Actor und Critic streng von einander getrennt und teilen keine gemeinsamen Parameter. Gelegentlich werden den beiden Netzen (Actor bzw. Critic) noch ein weiteres Netz vorgeschoben. In diesem Fall können Actor und Critic gemeinsame Parameter besitzen.
Das Actor- bzw. Policy-Network ist für die Bestimmung der Policy zuständig. Anders als bei Value-based RL-Verfahren wird diese direkt bestimmt und kann auch direkt angepasst werden. Die Policy wird als eine Wahrscheinlichkeitsverteiling über alle möglichen Actions vom Actor-NN zurückgegeben \cite{DRL}.\\
Das Critic- bzw. Value-Network evaluiert die Actions, welche vom Actor-Network bestimmt worden sind. Genauer gesagt, schätzt das Value-Network die sogenannte ''Discounted Sum of Rewards'' zu einem Zeitpunkt $t$, basierend auf dem momentanen State $s$, der dem Value-Network als Input dient. ''Discounted Sum of Rewards'' wird im späteren Verlauf noch weiter vorgestellt und erklärt.

\subsection{PPO Training Objective Function}
Nun da einige Grundlagen näher beleuchtet worden sind, ist das nächste Ziel die dem PPO zugrunde liegende mathematische Funktion zu verstehen, um im späteren eine eigene Implementierung des PPO durchführen zu können und um einen objektiveren Vergleich der zwei RL-Verfahren durchführen zukönnen.
\\Der PPO basiert auf den folgenden matematischen Formel, welche den Loss eines Updates bestimmt:
\begin{align}
	\label{PPO_Training_Loss}
	L^\text{PPO}_{t} (\theta) = L^\text{CLIP + VF + S}_{t} (\theta) = \mathbb{\hat{E}}_{t} [L^{\text{CLIP}}_{t}(\theta) - c_{1}L^{\text{VF}}_{t} + c_{2}S[\pi_{\theta}](s_{t})]
\end{align}
Dabei besteht die Loss-Function aus drei unterschiedlichen Teilen. Zum einen aus dem Actor-Loss bzw. Policy-Loss bzw. Main Objective Function $L^{\text{CLIP}}_{t}(\theta)$, zum anderen aus dem Critic-Loss bzw. Value-Loss $L^{\text{VF}}_{t}$ und aus dem Entropy Bonus $S[\pi_{\theta}](s_{t})$. Die Main Objective Function sei dabei durch folgenden Term gegeben \cite{PPO}.
\begin{align}
	L^\text{CLIP}_{t} (\theta) = \mathbb{\hat{E}}_{t} [ \min(r_{t}(\theta) \hat{A}_{t}(s, a), \text{clip}(r_{t}(\theta), 1 - \epsilon, 1 + \epsilon) \hat{A}_{t}(s, a))]
\end{align}

\subsubsection{Formelelemente}
Um die dem PPO zugrundeliegende Update Methode besser zu verstehen folge eine Erklärung ihrer einzelnen mathematischen Elemente.
Die einzelnen Erklärungen basieren auf den PPO Paper \cite{PPO}.
\begin{longtable}[h]{|p{4cm}|p{\linewidth - 5cm}|}
	\caption{Formelelemente}
	\label{tab:Formelelemente} 
	\endfirsthead
	\endhead
	\hline
	Symbol & Erklärung \\
	\hline
	$\theta$ & Theta beschreibt die Paramenter aus denen sich die Policy des PPO ergibt. Sie sind die Gewichte, welche das Policy-NN definiert. \\
	\hline
	$\pi_{\theta}$ & Die Policy bzw. Entscheidungsfindungsregeln sind eine Wahrscheinlichkeitsverteilung über alle möglichen Actions. \cite{Sutton1998} \\
	\hline
	$L^\text{CLIP} (\theta)$ & $L^\text{CLIP} (\theta)$ bezeichnetz den sogenannten Policy Loss, welche in Abhängigkeit zu der Policy $\pi_{\theta}$ steht. Dabei handelt es sich um einen Zahlenwert, welcher den Fehler über alle Parameter approzimiert. Dieser wird für das Lernen des Netzes benötigt. \\
	\hline
	$t$ & Zeitpunkt \\
	\hline
	$\mathbb{\hat{E}}_{t}$ & $\mathbb{\hat{E}}_{t}$ gibt an, dass es sich um einen Erwartung handelt, welche einem empirischen Durchschnitt zugrundeliegt. \\
	\hline
	$r_{t}(\theta)$ & Quote zwischen alter Policy (nicht als Abhängigkeit angegeben, da sie nicht verändert werden kann) und aktueller Policy zum Zeitpunkt $t$. Daher auch Probability Ratio genannt. \\
	\hline
	$\hat{A}_{t}(s, a)$ & erwarteter Vorteil bzw. Nachteil einer Action $a$, welche im State $s$ ausgeführt wurde. welcher sich in Abhängigkeit von dem State $s$ und der Action $a$ befindet. \\
	\hline
	$\text{clip}$ & Mathematische Funktion zum Beschneidung eines Eingabewertes. Clip setzt eine Ober- und Untergrenze fest. Sollte ein Wert der dieser Funktion übergeben wird sich nicht mehr in diesen Grenzen befinden, so wird der jeweilige Grenzwert zurückgegeben. \\
	\hline
	$\epsilon$ & Epsilon ist ein Hyperparameter, welcher die Ober- und Untergrenze der Clip Funktion festlegt. Gewöhnlich wird für $\epsilon$ ein Wert zwischen $0.1$ und $0.2$ gewählt. \\
	\hline
	$\gamma$ & Gamma bzw. Abzinsfaktor ist ein Hyperparameter, der die Zeitpräferenz des Agenten kontrolliert. Gewöhnlich liegt Gamma $\gamma$ zwischen 0.9 bis 0.99. Große Werte sorgen für ein weitsichtiges Lernen des Agenten wohingegen ein kleine Werte zu einem kurzfristigen Lernen führen \cite{Sutton1998}. \\
	\hline
\end{longtable}

\subsubsection{Advantage}
Der erste behandelt den Advantage $\hat{A}_{t}(s, a)$. Dieser wird durch die Subtraktion der Discounted Sum of Rewards bzw. des Return $R_{t}$ und dem Baseline Estimate $b(s_{t})$ bzw. der Value function berechnet \cite{DBLP:journals/corr/MnihBMGLHSK16, PPO}. 
\begin{align}
	\hat{A}_{t}(s, a) = R_{t} - b(s_{t})
\end{align}
Der Advantage gibt ein Mass an, um wie viel besser oder schlechter eine Action war, basierend auf der Erwartung der Value function bzw. des Critics. Es wird also die Frage beantwortet, ob eine gewählte Action $a$ im State $s_{t}$ zum Zeitpunkt $t$ besser oder schlechter als erwartet war \cite{DBLP:journals/corr/MnihBMGLHSK16}.

\subsubsection{Return} 
\label{sec:Return}
Der Return $R_{t}$ stellt dabei die Summe der Rewards in der gesamten Spieleepisode von dem Zeitpunkt $t$ an dar. Diese kann ermittelt werden, da alle Rewards, durch das Sammeln von Daten, bereits bekannt sind. Des Weiteren werden die einzelnen Summanden mit einem Discount Factor $\gamma$ multipliziert, um die Zeitpräferenz des Agenten besser zu steuern. Gamma liegt dabei gewöhnlich zwischen einem Wert von 0.9 bis 0.99. Kleine Werte für Gamma sorgen dafür, dass der Agent langfristig eher dazu tendiert Actions bzw. Aktionen zu wählen, welche unmittelbar zu positiven Reward führen. Entsprechend verhält es sich mit großen Werten für Gamma \cite{Sutton1998}.

\subsubsection{Baseline Estimate} 
\label{sec:Baseline_Estimate}
Der Baseline Estimate $b(s_{t})$ oder auch die Value function ist eine Funktion, welche durch ein NN realisiert wird. Es handelt sich dabei um den Critic der Actor-Critic-Verfahrens. Die Value function versucht eine Schätzung des zu erwartenden Discounted Rewards $R_{t}$ bzw. des Returns, vom aktuellen State $s_{t}$, zu bestimmen. Da es sich hierbei um die Ausgabe eines NN handelt, wird in der Ausgaben immer eine Varianz bzw. Rauschen vorhanden sein \cite{DBLP:journals/corr/MnihBMGLHSK16}.

\subsubsection{Probability Ratio}
Die Probability Ratio $r_{t}(\theta)$ ist der nächste Baustein zur Vervollständigung der PPO Main Objective Function. 
\begin{align}
	\label{for:Probability_Ratio}
	r_{t}(\theta) = \frac{\pi_{\theta}(a_{t}|s_{t})}{\pi_{\theta_{\text{old}}}(a_{t}|s_{t})}
\end{align}
In normalen Policy Gradient Methoden bestehe ein Problem zwischen der effizienten Datennutzung und dem Updaten der Policy. Dieses Problem tritt z.B. im Zusammenhang mit dem Advantage Actor Critic (A2C) Aglorithmus auf und reklementiert das effiziente Sammeln von Daten. So ist des dem A2C nur möglich von Daten zum lernen, welche on-policy (unter der momentanen Policy) erzeugt wurden. Das Verwenden von Daten, welche unter einer älteren aber dennoch ähnlichen Policy gesammelt wurden, ist daher nicht zu empfehlen.\\
Der PPO bedient sich jedoch eines Tricks der Statistik, dem Importance-Sampling (IS, deutsch: Stichprobenentnahme nach Wichtigkeit). 
Unter der Zuhilfenahme des IS ist es dem PPO nun möglich, auch Daten zu nutzen, welche unter ähnlichen älteren Policy gesammelt wurden. Dies eröffnet die Möglichkeit der Implementierung eines Replay-Buffers bzw. Memorys. Des Weiteren können nun generierte Daten mehrfach für Updates der Policy genutzt werden, was die Menge an Daten, welche nötig ist, um ein gewisses Ergebnis zu erreichen, stark minimiert. Der PPO Algorithmus ist durch die Umstellung auf die Probability Ratio effizienter in der Nutzung von Daten geworden.

\subsubsection{Surrogate Objectives}
Der Name Surrogate (Ersatz) Objective Function ergibt sich aus der Tatsache, dass die Policy Gradient Objective Function des PPO nicht mit der logarithmierten Policy $\log_{\pi_{\theta}}(a_{t}|s_{t})\hat{A}_{t}$ arbeitet, wie es die normale Policy Gradient Methode vorsieht, sondern mit dem Surrogate der Probability Ratio $r_{t}(\theta)$ \ref{for:Probability_Ratio} zwischen alter und neuer Policy.\\
Intern beruht der PPO auf zwei sehr ähnlichen Objective Functions, wobei die erste dieser beiden
\begin{align}
	\label{func:surrogate_1}
	r_{t}(\theta) \hat{A}_{t}(s, a)
\end{align}
der normalen TRPO Objective Function entspricht, ohne die, durch den TRPO vorgesehende, KL-Penalty \cite{TRPO, PPO}. 
Die alleinige Nutzung dieser Objective Function hätte jedoch destruktiv große Policy Updates zufolge. Aus diesem Grund haben John Schulman et al. eine zweite Surrogate Objective Function, dem PPO Verfahren hinzugefügt. 
\begin{align}
	\text{clip}(r_{t}(\theta), 1 - \epsilon, 1 + \epsilon) \hat{A}_{t}(s, a)
\end{align}
Die einzige Veränderung im Vergleich zur ersten Objective Function \ref{func:surrogate_1} ist, dass eine Abgrenzung durch die Clip-Funktion eintritt. Sollte sich die Probability Ratio zuweit von 1 entfernen, so wird $r_{t}(\theta)$ entsprechende auf $1 - \epsilon \text{ bzw. } 1 + \epsilon$ begrenzt. Das hat zufolge, dass der Loss $L^{\text{CLIP}}_{t}(\theta)$ ebenfalls begrenzt wird, sodass es zu keinen großen Policy Updates kommen kann. Es wird sich daher in einer Trust Region bewegt, da man sich nie allzu weit von der ursprünglichen Policy entfernt \cite{TRPO, PPO}.

\subsubsection{Zusammenfassung der PPO Training Objective Function}
$L^\text{CLIP} (\theta) = \mathbb{\hat{E}}_{t} [ \min(r_{t}(\theta) \hat{A}_{t}(s, a), \text{clip}(r_{t}(\theta), 1 - \epsilon, 1 + \epsilon) \hat{A}_{t}(s, a))]$
Insgesamt ist der Actor-Loss ein Erwartungswert, welcher sich empirisch über eine Menge an gesammelten Daten ergibt. Dies wird durch $\mathbb{\hat{E}}_{t}$ impleziert. Dieser setzt sich aus vielen einzelnen Losses zusammen. Diese einzelnen Losses  sind das Minimum der beiden Surrogate Objective Functions zusammen. Dies sorgt dafür, dass keine zu großen Losses die Policy zuweit von dem Entstehungspunkt der Daten wegführen (Trust Region) \cite{PPO}.\\
\\Des Weiteren wird für den PPO Training Loss auch noch der Value-Loss benötigt. Dieser setzt sich folgendermaßen zusammen:
\begin{align}
	L^{\text{VF}}_{t} = (V_{\theta}(s_{t})-V_{t}^{targ})^2 \text{ wobei } V_{t}^{targ} = r_{t}(\theta)
\end{align} 
Der letzte Teil, welcher für die Bestimmung des PPO Training Loss benötigt wird, ist der Entropy Bonus. Dabei handelt es sich um die Entropie der Policy. Diese lässt sich normal bestimmen.\\
Es ergibt sich damit der bereits oben erwähnte \ref{PPO_Training_Loss} PPO Training Loss \cite{PPO}:
\begin{center}
	$L^\text{PPO}_{t} (\theta) = L^\text{CLIP + VF + S}_{t} (\theta) = \mathbb{\hat{E}}_{t} [L^{\text{CLIP}}_{t}(\theta) - c_{1}L^{\text{VF}}_{t} + c_{2}S[\pi_{\theta}](s_{t})]$
\end{center}

\section{Deep Q-Learning}
Das Q-Learning ist eine weitere Art des RL, welches eine gesamte Familie an RL-Verfahrewn umfasst. 
Namensgebend für die Familie des Q-Learnings sind jedoch die sogenannten Q-Values, welche die Aktionswert einer Action $a$ im State $s$ angeben. Deutlich wichtiger im weiteren Verlauf des Deep Q-Learnings sind jedoch die Zustandswerte, welche von der Zustandswertefunktion $V(s)$ ermittelt werden. Bei ihr handelt es sich, bezogen auf das Deep Q-Learning um eine NN, dass im Verlauf triniert wird.\\
Um das Q-Learning auch hier vollständig zu verstehen und um einen Eindruck des zu erwartenden Implementierungsaufwandes zu erhalten, ist es nötig die dem DQN zugrundeliegenden Formeln zu verstehen.
Angestammtes Ziel des Q-Learnings ist es für alle State-Action-Paare $(s,a)$ ein $Q$ zu ermitteln \cite{DRL}. Daher gilt:
\begin{align}
	\label{Q-Value_Formel}
	Q(s,a) = \mathbb{E}_{s'\sim s}[r(s,a) + \gamma V(s')] = \sum_{s' \in S}^{}p_{a,s \to s'}(r(s,a) + \gamma V(s'))
\end{align}
Der Aktionswert für den State $s$ und die Action $a$ setzt sich auf dem zu erwartenden unmittelbaren Reward $r(s,a)$ und der diskontierten (abgezinsten) langfristige Reward des Folgezustandes $s'$ zusammen. Dabei ist $\gamma$ der Abzinsfaktor und der langfristige Reward des Folgezustand wird durch $V(s')$ dargestellt.\\
Diese Aktionswerte $Q(s,a)$ können nun herangezogen werden um das Value-NN zu trainieren, da der Value $V$ eines State $s$ durch den maximalen Aktionswert, über alle Actions, definiert ist \cite{DRL}.
\begin{align}
	V(s) = \max_{\substack{a \in A}}Q(s,a)
\end{align}
Der Value eines States ist daher der maximale Aktionswert desselbigen. Das bedeutet, dass der Aktionswert eines beliebigen State dem Wert entspricht, welcher sich mit der besten Action des States erzielen lässt.\\
Problematisch ist jedoch die Bestimmung der Q-Values, da sich \ref{Q-Value_Formel} schlecht in High-level-Languages wie Python, Java oder C++ implementieren lässt. In der Praxis wird daher ein rekursives Vorgehen dazu verwendet, um die Q-Values, über den Trainingsverlauf, zu erlernen \cite{DRL}.
 \begin{align}
 	\label{Q-Value_Formel_r}
 	Q(s,a) = r(s,a) + \gamma \max_{\substack{a \in A}}Q(s',a')
 \end{align}
Bei der letzen Formel handelt es sich auch um die durchzuführende Prozedue, um das Value-NN zu trainieren \cite{DRL}.
Da es sich beim Q-Learning um value-based RL-Verfahren \ref{RL_policy_value} handelt, muss kein weiteres NN aufgesetzt werden, wie es beispielsweise beim PPO der Fall ist. Im Vergleich zum voher vorgestellten PPO ist der DQL Algorithmus ein deutlich simpleres Verfahren. Die Tatsache, dass es sich bei dem DQL auch noch um eine off-policy Verfahren handelt, lässt unter anderem auch das Lernen von älteren Daten zu. Im Vergleich zum PPO ist der DQL-Agent dazu jedoch, in einem deutlich größeren Maß, imstande. Auch sehr weit zurückliegende Daten können vom DQL-Algorithmus zum Lernen verwendet werden. Dies bedingt einige Details, die in der Implmentierung zum Tragen kommen.\\





\chapter{Anforderungen} \label{chap:Anforderungen}
In \autoref{chap:Grundlagen} wurden Grundlagen für die weiteren Vergleiche der Reinforcement Learning Agenten gelegt, welche auf zwei unterschiedlichen Algorithmen (DQN und PPO) basieren.\\
Um diese Vergleiche zu realisieren, soll ein System entwickelt werden, dass diese durchführen, festhalten und auswerten kann. Dieses soll aus einem Environment, mehreren Agenten beider Algorithmen sowie aus statistischen Analysekomponenten zur Leistungsbestimmung bestehen. Zuzüglich sollen weitere Anforderungen an die Evaluation gestellt werden, um die Vergleichbarkeit sicherzustellen.

\section{Anforderungen an das Environment} \label{sec:Anforderungen_Env}
In diesem Abschnitt werden die Anforderungen an das Env dargestellt. Neben der Hauptanforderung, dass das Spiel Snake implementiert werden soll, ergeben sich weitere zusätzliche Anforderungen.

\subsection{Standardisierte Schnittstelle} \label{subsec:Anforderungen_Schnittstelle}
Das Env soll eine standardisierte Schnittstelle besitzen, sodass drei Kommunikationskanäle implementiert werden \fullref{subsubsec:Grundlagen_Environment}. Es soll in der Lage sein, Aktionen zu empfangen. Des Weiteren soll eine Observation und ein Reward an den Agenten übergeben werden. Diese Standardisierung erleichtert die Verwendbarkeit, auch bei anderen Algorithmen.

\subsection{Funktionalitäten} \label{subsec:Anforderungen_Funktionalität_Env}
Das Env soll die folgenden Funktionalitäten implementieren.

\subsubsection{Aktionsausführung} \label{subsubsec:Anforderungen_Aktionsausführung}
Das Env muss eine Funktionalität beinhalten, die eine Aktion ausführen kann. Diese Aktionsausführung muss sich nach den Regeln des Spiels Snake richten \fullref{sec:Grundlagen_Game_of_Snake}.

\subsubsection{Reset} \label{subsubsec:Anforderungen_Reset}
Das Env muss eine Reset Funktionalität implementieren, um einen erbrachten Spielfortschritt zurückzusetzen. Dies ist für den Spielablauf unentbehrlich.

\subsubsection{Render} \label{subsubsec:Anforderungen_Render}
Das Env muss eine Render Funktionalität implementieren, um eine Visualisierung des Spiels Snake zu ermöglichen. Diese dient der besseren Evaluation und Demonstration.

\section{Anforderungen an die Agenten} \label{sec:Anforderungen_Agenten}
In diesem Abschnitt werden die Anforderungen an die Agenten, welche auf dem PPO bzw. DQN Algorithmus basieren, dargestellt.

\subsection{Funktionalitäten} \label{subsec:Anforderungen_Funktionalitäten_Agent}
Die Agenten müssen folgende Funktionalitäten implementieren.

\subsubsection{Aktionsbestimmung} \label{subsubsec:Anforderungen_Aktionsbestimmung}
Die Agenten müssen in der Lage sein, aus einer Observation eine Aktion zu bestimmen, welche wiederum dem Env übergeben werden muss, um einen Spielfortschritt erzielen zu können.

\subsubsection{Lernen} \label{subsubsec:Anforderungen_Lernen}
Die Agenten müssen fähig sein, auf Grundlage vergangener Spielepisoden zu lernen und damit ihre Spielergebnisse zu verbessern.

\subsubsection{Parametrisierung} \label{subsubsec:Anforderungen_Parametrisierung}
Das System muss die Möglichkeit besitzen, mehrere Agenten des gleichen Algorithmus zu erstellen, welche sich jedoch durch ihre verwendeten Hyperparameter unterscheiden. Diese Definition von Agenten ist in der Evaluation zu berücksichtigen und dient damit einer besseren Vergleichbarkeit.

\subsection{Diversität der RL Algorithmen} \label{subsec:Anforderungen_Diversität}
Um nicht nur Agenten eines Algorithmus untereinander zu vergleichen, sondern auch den Vergleich zu anderen Algorithmen zu erbringen, sollen ein DQN und PPO Algorithmus miteinander verglichen werden. 
Diese bieten sich, wie in \autoref{sec:Grundlagen_PPO} und \autoref{sec:Grundlagen_Deep_Q_Learning} beschrieben, für den Vergleich an.

\section{Anforderungen an die Datenerhebung} \label{sec:Anforderungen_an_die_Datenerhebung}
In diesem Teil sollen Anforderungen an die statistische Datenerhebung und an die damit verbundenen Analysekomponenten gestellt werden.

\subsection{Mehrfache Datenerhebung} \label{subsec:Anforderungen_mehrfache_Datenerhebung}
Die Datenermittlung muss für jeden einzelnen Agenten mehrfach durchgeführt werden, um die Validität der Messung zu gewährleisten. \citep[S. 135]{DL}

\subsection{Datenspeicherung} \label{subsec:Anforderungen_Datenspeicherung}
Damit aus den erzeugten Test- und Trainingsdaten statistische Schlüsse gezogen werden können, ist es wichtig, dass diese gespeichert werden. Da jedoch die Menge an Daten schnell riesige Dimensionen annehmen würde, sollen stellvertretend nur die Daten ganzer Spiele gespeichert werden. Diese Strategie stellt einen Kompromiss zwischen Vollständigkeit und effizientem Speicherplatzmanagement dar.\\
Das System soll folgende Daten speichern:
\begin{longtable}[h]{|p{4cm}|p{\linewidth - 5cm}|}
	\hline
	Daten & Erklärung \\
	\hline
	steps & Die in einem Spiel durchgeführten Züge. Diese geben in der Evaluation später Aufschluss über die Effizienz und weisen auf Lernfehler der Agenten hin, wie beispielsweise das Laufen im Kreis.\\
	\hline
	apples & Die Anzahl der gefressenen Äpfel in einem Spiel ist ein maßgeblicher Evaluationsfaktor zur Einschätzung des Lernerfolges.\\
	\hline
	wins & Hat der Agent das Spiel gewonnen? Dieser Wert stellt die Endkontrolle des Agenten dar. Er gibt Aufschluss über das Konvergenzverhalten.\\
	\hline
	\caption[Darstellung der zu erhebende Daten]{Darstellung der zu erhebenden Daten für die statistische Auswertung}
	\label{tab:Anforderungen_Datenerhebung} 
\end{longtable}

\section{Anforderungen an die Evaluation} \label{sec:Anforderungen_an_die_Evaluation}
Bei der Evaluation soll der optimale Agent für jedes Evaluationskriterium ermittelt und durch aussagekräftige Statistiken untermauert werden.
Das Kriterium der Performance wurde von der verwendeten Literatur \cite{Autonomous_Agents_in_Snake_Game_via_DRL} und \cite{UAV} übernommen, die anderen ergeben sich aus den zu speichernden Daten.
Die einzelnen Kriterien lauten:
\begin{longtable}[h]{|p{4cm}|p{\linewidth - 5cm}|}
	\hline
	Kriterium & Erläuterung \\
	\hline
	Performance & Welcher Agent erreicht das beste Ergebnis? Im Sachzusammenhang mit dem Spiel Snake bedeutet dies: Welcher Agent frisst die meisten Äpfel, nachdem er trainiert wurde?\\
	\hline
	Siegrate & Welcher Agent schafft die Spiele mit einem Sieg zu beenden? Im Gegensatz zur Performance, gibt die Siegrate Aufschluss über das Konvergenzverhalten.\\
	\hline
	Robustheit & Welcher Agent kann in einer modifizierten Umgebung die größte Performance erreichen? In Bezug auf Snake bedeutet dies: Welcher Agent ist in der Lage, auf einem größeren bzw. kleineren Spielfeld die meisten Äpfel zu fressen? Dies ist bei Real World Applikationen ein wichtiger Faktor, da sich unbemannte Drohnen auch in unbekannten Umgebungen zurechtfinden müssen.\\
	\hline
	Effizienz & Welcher Agent löst das Spiel mit der größten Effizienz? Bezogen auf das Spiel Snake bedeutet dies: Welcher Agent ist in der Lage, die Äpfel mit möglichst wenig Schritten zu fressen? Dieser Wert ist besonders in Real World Applikationen von Interesse, da beispielsweise selbstfahrende Autos ihre Ziele in einer möglichst geringen Strecke erreichen sollen, um Energie und Zeit zu sparen.\\
	\hline
	\caption{Auflistung der Evaluationskriterien}
	\label{tab:Anforderungen_Kriterien} 
\end{longtable}
\chapter{Agenten}
Ein zentraler Aspekt der eines Vergleiches von verschiedenen RL-Agenten ist die genaue Definition der einzelnen Agenten. Basierende auf den Grundlagen \ref{sec:Agent} soll in diesem Kapitel der Begriff vervollständigt und die zu vergleichenden Agenten sollen vorgestellt werden.\\
Erste statistische Erhebungen haben gezeigt, dass die ausgewählten Hyperparameter einen immensen Einfluss auf das Verhalten der Agenten haben. Bestätigt wird diese Aussage durch die Quelle \cite{Sutton1998}. Ein Vergleich zwischen DQN und PPO mit wahllos gewählten Hyperparametern ist folglich wenig aussagekräftig. Daher ist auch die Definition des Begriffs Agent, welcher nur zwischen DQN und PPO diffenenziert, unzureichend.\\
\\Angebracht wäre eine neuer erweiterte Definition des Begriffs Agent für diese Ausarbeitung. Diese soll um den entscheidenen Faktor der Hyperparameter erweitert werden. Ein Agent wird daher nicht mehr alleinig durch seine Art (Q-Learning oder Policy Gradient bzw. DQN oder PPO) definiert, sondern ebenfalls durch die ausgewählten Hyperparameter.Eine Analogie aus dem Tierreich sollte hier Klarheit verschaffen.\\
\\Im Tierreich gibt es Hunde und Katzen. Diese stellen die RL-Klassen, wie z.B. Q-Learning- oder Policy Gradient Verfahren, dar. Sieht man jedoch genauer hin, so unterscheiden sich die Hunde und Katzen durch ihre jeweiligen Rassen, wie z.B. Pudel und Dalmatiner bei den Hunden und Maine Coons und Norwegische Waldkatzen bei den Katzen. Diese stellen die Algorithmusklassen, wie z.B. DQN oder PPO, dar. Dennoch unterscheiden sich auch Hunde und Katzen der selben Rasse untereinander, nämlich in ihrer DNS. Diese stellt die letzte Differenzierungsebene der Agenten dar. Im Sachzusammenhang stellen die Hyperparameter und Attribute, wie beispielsweise die Netzstruktur, die DNS eines Agenten dar.\\
Soll nun also ein Vergleich zwischen verschiedenen Agenten vollzogen werden, so gilt es als erstes die einzelnen Agenten zu definieren, daher ihre RL-Klasse, Algorithmusklasse und Hyperparameter zu bestimmen.
 
%--------------------------
%---Literaturverzeichnis---
%--------------------------

%Dokument vollständig setzen:
%latex
%bibtex
%latex
%latex

\cleardoublepage
\phantomsection
\addcontentsline{toc}{chapter}{Literaturverzeichnis} %Literaturverzeichnis im Inhaltsverzeichnis aufführen

\bibliography{Bibliographie/Bibliographie} %Bibliographie-Datei ohne .bib-Endung angeben!

%------------
%---Anhang---
%------------

\appendix %nummerierte Kapitel im Anhang mit Buchstaben statt Zahlen

%\chapter{Einleitung}\label{sec:beispiel}
Das Maschine Learning ist weltweit auf dem Vormarsch und große Unternehmen investieren riesige Beträge, um mit KI basierten Lösungen größere Effizienz zu erreichen. Auch der Bereich des Reinforcement Learning gerät dabei immer mehr in das Blickfeld von der Weltöffentlichkeit. Besonders im Gaming Bereich hat das Reinforcement Learning schon beeindruckende Resultate erbringen können, wie z.B. die KI AlphaGO, welche den amtierenden Weltmeister Lee Sedol im Spiel GO besiegt hat \cite{UAV}. In Anlehnung an die vielen neuen Reinforcement Learning Möglichkeiten, die in der letzten Zeit entwickelt wurden und vor dem Hintergrund der immer größer werdenden Beliebtheit von KI basierten Lösungsstrategien, soll es in dieser Bachelorarbeit darum gehen, einzelnen Reinforcement Learning Agenten, mittels statistischer Methoden, genauer zu untersuchen und den optimalen Agenten für ein entsprechendes Problem zu bestimmen.

\section{Motivation} \label{sec:Motivation}
In den letzten Jahren erregte das Reinforcement Learning eine immer größere Aufmerksamkeit. Siege über die amtierenden Weltmeister in den Spielen GO oder Schach führten zu einer zumehmenden Beliebtheit des RL. Neue Verfahren, wie z.B. der Deep Q-Network Algorithmus auf dem Jahr 2015 \cite{DQN}, der Proximal Policy Optimization (PPO) aus dem Jahr 2017 \cite{PPO} oder der Soft Actor Critic (SAC) aus dem Jahr 2018 \cite{SAC}, haben ihr Übriges getan, um das RL auch in anderen Bereichen weiter anzusiedeln, wie z.B. der Finanzwelt, im autonomen Fahren, in der Steuerung von Robotern, in der Navigation im Web oder als Chatbot \cite{DRL_Lapan}.\\
Durch die jedoch große Menge an RL-Verfahren gerät man zunehmend in die problematische Situaion, sich für einen diskreten RL Ansatz zu entscheiden. Weiter erschwert wird dieser Auswahlprozess noch durch die Tatsache, dass die einzelnen Agenten jeweils untereinander große Unterschiede aufweisen. Auch existieren häufig mehrere Ausprägungen eines RL-Verfahrens, wie z.B. der Deep Q-Network Algorithmus (DQN) und der Double Deep Q-Network Algorithmus (DDQN). Die Wahl des passenden Agenten kann großen Einfluss auf die Performance und andere Bewertungskriterien haben, (\cite{Exploration_of_Reinforcement_Learning_to_SNAKE}), deshalb soll in dieser Ausarbeitung eine Vorgehen, welches auf einem Vergleich beruht, entwickelt werden, dass den optimalen Agenten für eine entsprechendes Problem bestimmt.\\
\\Eine potenzielle Umgebung, in welcher Agenten getestet und verglichen werden können, ist das Spiel Snake.
Mit der Wahl dieses Spieles ist zusätzlich zu dem oben erwähnten Mehrwert noch ein weiterer in Erscheinung getreten.\\ 
So interpretieren neue Forschungsansätze das Spiel Snake als ein dynamisches Pathfinding Problem. Mit dieser Art von Problemen sind auch unbemannte Drohne (UAV Drohnen) konfrontiert, welche beispielsweise Menschen in komplexen Katastrophensituationen, wie z.B. auf explodierten Rohölbohrplattformen oder Raffinerien, finden und retten sollen. Auch kann das Liefen von wichtigen Gütern, beispielsweise medizinischer Natur, in solche Gebiete kann durch die Forschung am Spiel Snake möglich gemacht werden.(\cite{UAV})\\
Somit kann der RL-Vergleich am Spiel Snake möglicherweise helfen, Menschen zu retten.

\section{Zielsetzung}
Ziel der Arbeit ist es für die Umgebung Sanke einen optimalen Agenten zu bestimmen, welcher spezifische Anforderungen erfüllen soll.
Damit ergib sich die folgende Forschungsfrage:\\
Wie kann an einem Beispiel des Spiels Snake für eine nicht-triviale gering-dimensionale Umgebung ein möglichst optimaler RL Agent ermittelt werden?\\
\\Basierend auf der Forschungsfrage ergibt sich ein Mehrwert für die Wissenschaft. Durch die Abnahme des Entscheidungsfindungsprozesses müssen Forscherinnen und Forscher wie auch Anwenderinnen und Anwender von RL-Verfahren nicht mehr unreflektiert irgendeinen RL Agenten auswählen, sondern können auf Grundlage der Methodik und der daraus hervorgehenden Daten den passenden Agenten bestimmen.

 %auskommentieren, falls nicht benötigt

%-------------
%---Schluss---
%-------------

\backmatter %Keine nummerierten Kapitel für z.B. Glossar, Index, Erklärung, etc.

\pagenumbering{gobble} %keine Seitenzahl auf der Erklärung
\chapter*{Erklärung}
Hiermit versichere ich, \autor{}, dass ich diese Arbeit selbständig verfasst und keine anderen als die angegebenen Quellen und Hilfsmittel benutzt habe. Außerdem versichere ich, dass ich die allgemeinen Prinzipien wissenschaftlicher Arbeit und Veröffentlichung, wie sie in den Leitlinien guter wissenschaftlicher Praxis der \universitaet{} festgelegt sind, befolgt habe.

\vspace{3cm}

{\def\arraystretch{2}
\begin{tabular}{p{8cm}}
\hline 
\autor \\ 
\end{tabular}
}

 %Erklärung auf der letzten Seite der gesamten Arbeit



\end{document}